\section{Detalles de Diseño}
\label{sec:detalles-diseno}

Este capítulo desarrolla los algoritmos que el capítulo~\ref{sec:solucion-propuesta}
presentó a alto nivel: el orden exacto de un \emph{frame}, la física y las
colisiones, el DSL de comportamiento, las reglas de combate, la carga de niveles y
la estrategia de testing.

\subsection{Pipeline de un frame}
\label{subsec:pipeline-frame}

Un \emph{frame} en \texttt{Playing} con \emph{delta time} activo es una sola
invocación de \texttt{advanceSimulationFrame}
(sección~\ref{subsec:solucion-nucleo}). El orden interno es fijo:

\begin{enumerate}
  \item \texttt{advanceFrame}: \texttt{runBehaviourStep} y \texttt{step}
    (plataformas especiales, física y colisiones)
  \item combate (contacto y \emph{melee} del \emph{player})
  \item proyectiles (lanzamiento del \emph{player} y \emph{enemy shot}s)
  \item \emph{falling hazard}s (peligros que caen desde el techo del nivel)
  \item transiciones de fase del jefe
  \item \texttt{resolvePickups} y victoria híbrida (\texttt{resolvePlayingWin})
  \item \emph{out-of-bounds} (caída fuera del nivel) y muerte o \emph{respawn}
  \item \texttt{resolveFramePhase} si victoria y pérdida de vida coinciden
\end{enumerate}

En fases \texttt{GameOver}, \texttt{LevelComplete} o \texttt{Victory}, o con
\emph{delta time} congelado, \texttt{updateGame} no invoca
\texttt{advanceSimulationFrame}. \texttt{resolveFramePhase} anula la victoria en
el mismo \emph{frame} ante \texttt{GameOver} o \emph{respawn}.
\texttt{updateGame} lee \texttt{GameConfig}, delega el \emph{frame} al dominio y escribe
el \texttt{GameState} con \texttt{modify}. No reexpone la pila \texttt{GameM}
(sección~\ref{subsec:solucion-monadas}). Las subsecciones siguientes detallan cada
bloque.

\subsection{Física cinemática}
\label{subsec:fisica}

La cinemática del \emph{player} y de los \emph{enemy} terrestres sigue la misma
ley en unidades del juego (px, px/s, px/s\textsuperscript{2}). Tras registrar si
el \emph{player} estaba en el suelo (\texttt{wasOnGround}), el motor aplica en
orden: \emph{input} horizontal ($v_x$ según signo del movimiento y velocidad máxima),
gravedad ($v_y' = v_y - g \cdot \Delta t$), salto (si hubo \emph{press} de salto y
\texttt{wasOnGround}, $v_y$ pasa al impulso autorizado) e integración de posición
($\mathit{pos}' = \mathit{pos} + \mathit{vel} \cdot \Delta t$). El salto solo
cuenta si la tecla se pulsa en ese \emph{frame}, no si permanece apretada.

La física anterior corre una vez por \emph{frame}. Solo la integración y la
colisión del \emph{player} se subdividen cuando $|v_y| \cdot \Delta t$ supera la
mitad de la altura del sólido más bajo. El número de sub-pasos se acota entre $1$
y $16$. Cada sub-paso usa una fracción de $\Delta t$. Así se reduce el \emph{túnel}
\emph{AABB} (atravesar una plataforma fina en un solo paso) en plataformas finas sin
repetir \emph{input} ni gravedad. Los \emph{enemy} voladores
integran posición sin gravedad ni colisión con plataformas. Los terrestres
reciben gravedad y la misma resolución \emph{AABB} que el jugador, sin sub-pasos extra.

\subsection{Colisiones AABB}
\label{subsec:colisiones}

Cada \emph{entity} usa una \emph{collision box} alineada a los ejes (\emph{AABB}). Las
plataformas se anclan en la esquina inferior izquierda (\emph{bottom-left}) y el
\emph{foot position} (punto de apoyo en el borde inferior de la caja) del jugador y de
los enemigos cae en el centro de ese borde.

Ante un solapamiento jugador--sólido, el motor elige el eje de \emph{menor
penetración}: empujar en $X$ para choques laterales y en $Y$ para suelo o techo.
Resolver siempre el eje vertical primero confundiría una pared con un aterrizaje
(o un techo con un apoyo). Cuando las penetraciones son casi iguales, un sesgo
favorece $Y$ para tratar el contacto como apoyo.

En el eje vertical interviene \texttt{vyBefore}, la componente vertical \emph{antes}
de integrar posición. Solo se aterriza o se corrige techo si \texttt{vyBefore}
$\leq 0$. Si el jugador subía ($\texttt{vyBefore} > 0$), un choque superior se
trata como golpe de techo, no como piso. Las plataformas se procesan ordenadas por
borde superior descendente. Se repite hasta ocho pasadas o hasta que no quede
solapamiento. Los enemigos terrestres comparten la regla. Los voladores la
ignoran.

\subsection{Plataformas especiales}
\label{subsec:plataformas-especiales}

\textbf{Plataformas móviles.} Cada una hace \emph{ping-pong} entre dos extremos con ancla
\emph{bottom-left} y velocidad constante. Se avanzan antes de la física del
\emph{player}. Si el jugador está en el suelo y apoyado sobre el tramo superior
de una móvil, su posición se desplaza con la plataforma ese \emph{frame}
(\emph{displacement carry}: el jugador ``sube'' con la plataforma sin copiar su velocidad al
saltar). Al saltar no hereda su velocidad.

\textbf{Plataformas desmoronables.} Permanecen sólidas para el jugador hasta que
descansa encima. Entonces arranca un contador fijo de frames (15 en la
implementación actual) y siguen sólidas durante la cuenta. Al terminar, dejan de
colisionar con el jugador, caen verticalmente y se eliminan bajo la línea de
muerte. Los enemigos las tratan como sólidas hasta que desaparecen.

\textbf{Arena de jefe.} Campo opcional en la \emph{level definition}. Al cruzar
el umbral horizontal el \emph{run} marca que el combate con el jefe está activo. Mientras
haya jefe vivo, paredes virtuales laterales confinan al \emph{player} dentro de la región
y la cámara no sale de esos límites.
El detalle de combate del jefe se resume en la sección~\ref{subsec:combate}.

\subsection{DSL de comportamiento}
\label{subsec:dsl}

Cada \emph{enemy} activo lleva un \texttt{BehaviourProgram}: un \emph{free monad}
sobre el \emph{functor} \texttt{EntityAction} (cada nodo del \emph{AST} es una acción
posible). El interpreter \texttt{runBehaviourStep} consume \emph{una} instrucción por
\emph{frame} y puede leer el \emph{world} para \emph{sensing} (medir distancias al
\emph{player} o al spawn). Luego \texttt{step} integra física. Los bucles infinitos
(patrullas, \emph{FSM}s reactivas) se escriben con \texttt{fix} y recursión en el
\emph{AST}, no con mutación. Una \emph{FSM} (\emph{finite state machine}) alterna
estados según condiciones del \emph{world}.

\begin{table}[htbp]
  \centering
  \caption{Primitivas del DSL (un paso por \emph{frame} salvo \texttt{WaitFrames}).}
  \label{tab:dsl-primitivas}
  \begin{tabular}{@{}p{4.4cm}p{8.6cm}@{}}
    \textbf{Instrucción} & \textbf{Efecto en un \emph{behaviour step}} \\
    \hline
    \texttt{SetVelocity} & Fija \texttt{enemyVel}. \\
    \texttt{WaitFrames} & Mantiene la velocidad y decrementa el contador. \\
    \texttt{IfPlayerWithinRange} & Elige rama según distancia horizontal al \emph{player}. \\
    \texttt{IfNearSpawn} & Elige rama según proximidad al \emph{spawn anchor} (punto de aparición). \\
    \texttt{MoveTowardPlayer} & Velocidad horizontal hacia el \emph{player}. \\
    \texttt{MoveTowardSpawn} & Velocidad horizontal hacia el spawn anchor. \\
    \texttt{MoveTowardPlayer2D} & Velocidad en el plano hacia el \emph{player} (voladores). \\
    \texttt{MoveTowardSpawn2D} & Velocidad en el plano hacia el spawn anchor (voladores). \\
    \texttt{FacePlayer} & Orienta hacia el \emph{player} y fija velocidad cero. \\
    \texttt{SetFacingTowardPlayer} & Orienta hacia el \emph{player} sin cambiar velocidad. \\
    \texttt{Shoot} & Dispara si el cooldown lo permite. \\
  \end{tabular}
\end{table}

Los programas compuestos viven en \texttt{Domain.Logic.BehaviourCatalog}:

\begin{itemize}
  \item \textbf{Patrulla horizontal} (\texttt{patrolHorizontal}): alterna
    velocidad $\pm speed$ con esperas de $n$ frames por tramo. Típica del caracol.
  \item \textbf{FSM reactiva terrestre} (\texttt{reactiveFsm}): si el
    \emph{player} entra en rango, persigue. Si está cerca del spawn, mira al
    \emph{player}. Si no, vuelve al spawn. Los \emph{presets} \texttt{chase} y
    \texttt{guard} (plantillas de archetype) comparten el mismo programa. La diferencia está en los
    parámetros de movimiento y en el \texttt{BehaviourTuning} opcional.
  \item \textbf{FSM aérea} (\texttt{flyingReactiveFsm}): igual en espíritu, con
    chase omnidireccional y patrulla horizontal en el spawn para murciélagos.
  \item \textbf{Arquero} (\texttt{archerProgram}): en rango, mira y dispara. Fuera
    de rango, mira y espera.
\end{itemize}

Sin \emph{preset} ni \emph{hint}, cada \texttt{kind} usa el programa por defecto de su
clase (\texttt{snail} patrulla, \texttt{golem} FSM terrestre, \texttt{bat} FSM aérea,
\texttt{archer} disparo, \texttt{bossGolem} y \texttt{bossBat} catálogo de fases).

\begin{table}[htbp]
  \centering
  \caption{Parámetros de \texttt{BehaviourTuning} y dónde se aplican.}
  \label{tab:behaviour-tuning}
  \begin{tabular}{@{}p{2.4cm}p{4.2cm}p{5.4cm}@{}}
    \textbf{Campo} & \textbf{Efecto} & \textbf{Módulo} \\
    \hline
    \texttt{speed} &
    Escala velocidades del programa \emph{DSL}. &
    \texttt{BehaviourCatalog.applyTuning} \\
    \texttt{reach} &
    Escala rangos de detección y de disparo. &
    \texttt{BehaviourCatalog.applyTuning} \\
    \texttt{toughness} &
    Escala puntos de salud del enemigo. &
    \texttt{BuildWorld} al materializar el \emph{world} \\
  \end{tabular}
\end{table}

Al construir el \emph{world}, \texttt{programForEnemyDef} elige el programa según
precedencia: \texttt{behaviourPreset} explícito $>$ \emph{preset} y \emph{tuning}
resueltos desde \texttt{behaviourHint} $>$ \texttt{defaultProgramForKind}. Un \emph{preset}
incompatible con la clase de movimiento cae a \texttt{idleProgram} (quedarse quieto).
Los jefes ignoran \emph{presets} y \emph{hints}. Cargan fases desde el catálogo interno
(sección~\ref{subsec:combate}).

\subsection{Combate y proyectiles}
\label{subsec:combate}

Tras la física, \texttt{resolveCombat} aplica daño por contacto: cualquier
solapamiento \emph{collision box} jugador--enemigo daña, sin pisotón. Respeta
\emph{frames} de invencibilidad del \emph{player}. El \emph{melee} extiende la caja del
\emph{player} según su orientación durante la ventana de ataque. El alcance y el
daño vienen de \texttt{CombatParams}.

\texttt{resolveProjectiles} integra proyectiles del \emph{player} (lanzamiento con
\emph{cooldown} y arco balístico) y \emph{enemy shot}s generados por \texttt{Shoot} en el
\emph{DSL}. Colisionan con sólidos y con cuerpos. Al impactar aplican daño y desaparecen.
Los \emph{falling hazard}s caen a velocidad autorizada, dañan por solapamiento y,
con \texttt{loopDelay} opcional en el JSON, reactivan en el spawn tras una espera.

La victoria híbrida exige solapamiento con la \emph{exit zone}, puntaje mínimo de la
\emph{level definition} y ausencia de jefe vivo (las tres condiciones juntas).
\texttt{resolveHazardsAndDeath} comprueba caída bajo la línea de muerte (derivada de la
plataforma más baja) y reduce vidas o marca \texttt{GameOver}. El \emph{respawn} ubica al
jugador en el \emph{spawn point}, elimina sus proyectiles y desactiva el modo de \emph{boss arena}.

Cada nivel admite como máximo un jefe. A diferencia del resto de enemigos, no
recibe \texttt{behaviourPreset} en el JSON: un catálogo interno fija umbrales de
salud y un programa por fase. Tras el combate del \emph{frame},
\texttt{resolveBossPhases} compara la salud con umbrales del catálogo y con un snapshot
pre-combate (\texttt{wBefore}). Si corresponde, sustituye el programa activo del jefe.

\subsection{Niveles}
\label{subsec:niveles}

La \emph{level definition} es JSON: \texttt{minScore}, \texttt{spawn}, listas de
plataformas (estáticas, móviles, desmoronables), \emph{enemy}s con \texttt{kind},
\texttt{behaviourPreset} o \texttt{behaviourHint} opcionales, \emph{pickup}s,
\emph{falling hazard}s, \texttt{exit} y \texttt{bossArena} opcional. Los
rectángulos usan \texttt{pos} (esquina inferior izquierda), \texttt{width} y \texttt{height}.

\textbf{Catálogo del \emph{run}.} Al arrancar, \texttt{BootstrapRun} arma la lista
de definiciones del \emph{run}: carga los JSON de referencia, puede sustituir
cada \emph{slot} por una definición generada vía API Anthropic (tres perfiles: intro,
desafío, jefe) y resuelve los \texttt{behaviourHint} antes de cualquier
\emph{frame}. Si la generación o la resolución fallan, el motor cae al JSON
fijo (\emph{fallback}) o al \emph{default} del \texttt{kind} sin abortar la partida.
Los niveles generados por API omiten \texttt{behaviourHint}. Cada enemigo queda con el
\emph{default} de su \texttt{kind}.

\textbf{\emph{Hints} en lenguaje natural.} El autor puede describir la IA de un
\emph{enemy} con \texttt{behaviourHint}. El adaptador traduce el texto a un
archetype (\texttt{patrol}, \texttt{chase} o \texttt{guard}) y a un
\texttt{BehaviourTuning}. \texttt{archer} es una \texttt{kind} con programa propio.
La construcción pura del \emph{world} solo ve \emph{preset} y \emph{tuning} ya
materializados en memoria.

El caso de uso \texttt{loadLevelFromText} decodifica con Aeson y llama a
\texttt{buildWorld}, que valida y materializa un \texttt{World} o devuelve
\texttt{LevelBuildError}. Errores de \emph{decode} o de construcción se elevan a
\texttt{GameError} en el borde con \texttt{IO}. El núcleo solo ve
\texttt{Text}.

Validaciones relevantes: \texttt{minScore} $\geq 0$, identificadores únicos por
colección, como máximo un \emph{enemy} de tipo jefe. Si hay \texttt{bossArena},
debe existir jefe en el nivel. Los jefes no admiten \texttt{behaviourPreset} ni
\texttt{behaviourHint}.

Fragmento con \emph{hint} ilustrativo:

\begin{lstlisting}[language={}]
{
  "minScore": 300,
  "spawn": {"x": -200, "y": 24},
  "platforms": [
    {"pos": {"x": -280, "y": 0}, "width": 405, "height": 8}
  ],
  "enemies": [
    {
      "id": 1,
      "kind": "snail",
      "pos": {"x": -120, "y": 8},
      "behaviourHint": "deambula lento y pesado, ajeno a quien pase"
    }
  ],
  "pickups": [],
  "exit": {"pos": {"x": 2540, "y": 8}, "width": 36, "height": 64}
}
\end{lstlisting}

\subsection{Tests}
\label{subsec:pruebas}

La suite usa \texttt{tasty} y \texttt{tasty-hunit}. Los módulos \texttt{*Test}
registran casos con descubrimiento automático en \texttt{test/Main.hs}. El
dominio se prueba con \emph{world}s fijos en \texttt{Domain.Fixtures} y
comparación \texttt{(@?=)}. La orchestration usa \texttt{runGameM} sin \texttt{IO}.
\texttt{updateGame} es polimórfica en \texttt{MonadReader} y
\texttt{MonadState}, lo que permite instancias mínimas en
\texttt{UpdateGameTest}.

No hay regresión visual ni cobertura exhaustiva: son \emph{smoke checks} (tests
mínimas de sanidad) sobre lógica delicada. Otros módulos cubren arena de jefe,
plataformas especiales y flujo de nivel con el mismo criterio. Ejemplos representativos:

\begin{table}[htbp]
  \centering
  \caption{Tests ilustrativos del motor.}
  \label{tab:pruebas}
  \begin{tabular}{@{}p{3.4cm}p{9.6cm}@{}}
    \textbf{Módulo} & \textbf{Qué verifica} \\
    \hline
    \texttt{StepTest} &
    \texttt{dt = 0} es identidad, aterrizaje, techo, pared y sub-pasos. \\
    \texttt{AabbTest} &
    Solapamiento y utilidades de caja. \\
    \texttt{BehaviourTest}, \texttt{EnemyFsmTest} &
    Un paso del interpreter y ramas de la FSM reactiva. \\
    \texttt{CombatTest}, \texttt{ProjectileTest} &
    Contacto, melee y disparos. \\
    \texttt{LevelLoadTest}, \texttt{LoadLevelTest} &
    JSON válido o rechazado y \texttt{buildWorld}. \\
    \texttt{FrameTest} &
    Orden del \emph{pipeline} y reglas de fase en un \emph{frame} completo. \\
    \texttt{BootstrapRunTest} &
    Fusión generado/\emph{fallback} y catálogo final del \emph{run}. \\
    \texttt{ResolveBehavioursTest} &
    Resolución de \texttt{behaviourHint} vía \emph{stub} del puerto. \\
    \texttt{BehaviourCatalogTest} &
    \texttt{programForEnemyDef} y \emph{tuning} sobre archetypes. \\
    \texttt{UpdateGameTest} &
    Simulación congelada en \texttt{GameOver}, \emph{pickup}s y \texttt{runFrames}. \\
  \end{tabular}
\end{table}

El test \texttt{unit\_stepZeroIsIdentityAtSpawn} fija el contrato de frame
congelado: con \emph{delta time} cero, \texttt{step} devuelve el \emph{world}
entrante sin importar el resto de la cadena.
